\documentclass[11pt,a4paper]{article}
\usepackage[T1]{fontenc}
\usepackage[utf8]{inputenc}
\usepackage{amsmath,amssymb,xcolor}
\usepackage[margin=25mm]{geometry}
\usepackage[colorlinks=true,urlcolor=blue]{hyperref}
\definecolor{nexorablue}{HTML}{2459B5}
\setlength{\parindent}{0pt}
\setlength{\parskip}{7pt}
\setlength{\emergencystretch}{2em}
\renewcommand{\familydefault}{\sfdefault}
\begin{document}

{\small\bfseries\color{nexorablue} NEXORA CPGE}\par
{\small MP / WKB-SPE-SI-01 / Énoncé}\par
{\LARGE\bfseries Réduire la boucle de retour unitaire}\par
\vspace{0.25em}\hrule\vspace{0.6em}

La chaîne directe a pour transfert $G(p)$. En retour unitaire négatif, la sortie $Y$ est prélevée et soustraite à la consigne $R$. Établir les deux fonctions $Y/R$ et $E/R$, où $E$ est l’écart en entrée de la chaîne directe.
\[ R\longrightarrow\boxed{\ R-Y\ }\overset{E}{\longrightarrow}\boxed{G(p)}\longrightarrow Y. \]

\bigskip
{\footnotesize Adaptation française par Nexora CPGE. Source : Contributeurs de Wikibooks,
\href{https://en.wikibooks.org/wiki/Control_Systems/Feedback_Loops\#Feedback_Loop_Transfer_Function}{\textit{Control Systems}}, Feedback Loops, dérivation du transfert de boucle.
\href{https://creativecommons.org/licenses/by-sa/4.0/}{CC BY-SA 4.0}.}
\end{document}
